\documentclass[conference]{IEEEtran}

\usepackage[utf8]{inputenc}
\usepackage[T1]{fontenc}
\usepackage{amsmath,amssymb}
\usepackage{graphicx}
\usepackage{booktabs}
\usepackage{tabularx}
\usepackage{hyperref}
\usepackage{cite}
\usepackage{balance}
\usepackage{listings}
\usepackage{xcolor}
\usepackage{multirow}
\usepackage{siunitx}

\lstset{
  basicstyle=\footnotesize\ttfamily,
  breaklines=true,
  frame=single,
  captionpos=b,
}

\hypersetup{
  colorlinks=true,
  linkcolor=black,
  citecolor=black,
  urlcolor=blue,
}

\begin{document}

\title{Forest Link Protocol: A Multi-Protocol Adaptive Mesh Network for Reliable Data Transfer in Dense Forest Environments}

\author{
  \IEEEauthorblockN{Po Haoting\IEEEauthorrefmark{1},
  Ong Tun Siang\IEEEauthorrefmark{1},
  Kenny Leck\IEEEauthorrefmark{1},
  Chia Wei Sheng\IEEEauthorrefmark{1}}
  \IEEEauthorblockA{\IEEEauthorrefmark{1}Singapore Institute of Technology\\
  CSC2106: IoT Protocols and Networks --- Group GP01\\
  \{2401280, 2402091, 2403543, 2400953\}@sit.singaporetech.edu.sg}
}

\maketitle

% ===========================================================================
\begin{abstract}
% ===========================================================================
This paper presents the design and implementation of Forest Link Protocol (FLP), a novel application-layer mesh protocol for IoT sensor deployments in environments where internet connectivity is intermittent and signal attenuation is high. Unlike standard IoT solutions that rely on continuous cloud connectivity or specialized LPWAN hardware, FLP utilizes commodity ESP32-S3 microcontrollers with dual-transport communication---ESP-NOW for short-range, high-bandwidth mesh relay and LoRa for long-range control signaling---to establish a self-healing, leader-electing tree topology. The protocol organizes nodes into a dynamic tree rooted at any node with verified internet connectivity. Nodes without internet access operate in a delay-tolerant mode, buffering data locally until the network reforms. FLP employs a Selective Repeat ARQ with Forward Error Correction to deliver reliable multi-megabyte file transfers over multiple hops. Our implementation on the LILYGO T3-S3 platform demonstrates that a 3\,MB file can be transferred over three hops in under three minutes in dense forest conditions, meeting the target requirement with significant margin.
\end{abstract}

% ===========================================================================
\section{Introduction}
\label{sec:introduction}
% ===========================================================================

Deploying sensor networks in remote, ecologically sensitive environments such as dense forests presents a unique set of engineering challenges that standard TCP/IP network stacks are ill-equipped to handle. The combination of heavy signal attenuation from water-heavy foliage, intermittent backhaul connectivity, and the need for long-lived battery-powered nodes creates a regime where off-the-shelf IoT protocols fail in predictable and costly ways.

The Forest Link Protocol (FLP) addresses these challenges by designing a custom application-layer mesh protocol tailored specifically to the ``Deep Field'' environment. Rather than adapting existing protocols such as LoRaWAN or standard MQTT over TCP, FLP is architected from first principles to be self-healing, leader-electing, and delay-tolerant by default.

\subsection{Problem Statement}

The primary issue we identified is the ``Green Wall'' effect, where dense and water-heavy foliage severely attenuates 2.4\,GHz WiFi signals. Standard IoT deployments in these areas face a ``Bandwidth-Reach Paradox'': high-bandwidth protocols (WiFi) cannot penetrate the vegetation over long distances, while long-range protocols (LoRa) lack the throughput necessary to move rich media or large data logs.

This results in disconnected or brittle networks that fail to deliver critical environmental data when signal quality fluctuates. Specific technical challenges include:

\begin{itemize}
  \item \textbf{Signal Attenuation} --- 2.4\,GHz signals are heavily absorbed by foliage, creating hidden-node problems.
  \item \textbf{Transient Backhaul} --- The internet gateway is often a single point of failure (e.g., a solar-powered 4G modem) that may sleep or fail without warning.
  \item \textbf{The Bandwidth Gap} --- LoRaWAN duty cycles severely limit payload size, making transfer of rich sensor datasets impractical.
  \item \textbf{The Reliability Gap} --- Standard TCP/IP meshes suffer from management collapse in forests; if the WiFi link drops, the entire protocol state fails.
\end{itemize}

\subsection{Relevance and Importance}

This research is vital for the Smart Forestry and Environmental Conservation industries. As climate change increases the frequency of wildfires and ecological shifts, society requires real-time, high-fidelity data (images, vibration logs, acoustics) from ``Deep Field'' environments. Academia currently lacks a robust, hybrid model that can reliably transfer multi-megabyte files across lossy, non-line-of-sight (NLOS) forest links without the infrastructure costs of satellite or cellular arrays at every node.

\subsection{Objectives}

The overarching goal is to implement and validate the Forest Link Protocol, focusing on:

\begin{enumerate}
  \item \textbf{Adaptive Protocol Selection} --- Developing an intelligent system that dynamically chooses communication strategy based on real-time topology, hop count, and RSSI measurements.
  \item \textbf{Self-Healing Topology} --- Implementing an implicit leader election mechanism that allows any node with internet connectivity to assume the root role automatically.
  \item \textbf{Delay Tolerant Networking} --- Ensuring that nodes without internet access buffer data locally and flush it to the cloud seamlessly upon reconnection.
  \item \textbf{Efficient Data Transport} --- Utilizing Selective Repeat ARQ with FEC parity to minimize retransmission overhead over lossy links.
\end{enumerate}

% ===========================================================================
\section{Literature Review}
\label{sec:litreview}
% ===========================================================================

\subsection{Current State of Knowledge}

The field of IoT mesh networking in constrained and remote environments has seen significant development across several protocol families. Low-Power Wide-Area Networks (LPWANs) still represent the dominant paradigm for remote IoT deployments, with LoRaWAN emerging as the most widely adopted standard across the four major LPWAN candidates---LoRa, NB-IoT, Sigfox, and LTE-M---in terms of global network operators and country coverage~\cite{maurya2025lorawan}.

For short-range mesh deployments, IEEE~802.15.4-based protocols such as Zigbee and Thread are widely studied. Delay Tolerant Networking (DTN) has been researched since the early 2000s via the Bundle Protocol (RFC~5050 and successors), which still remains the backbone of DTN, providing the store-carry-forward abstraction. Recent work emphasizes BPv7 (updated Bundle Protocol) for improved efficiency, extensibility, and security~\cite{elewaily2024dtn}.

\subsection{Significant Research Findings}

LoRaWAN has become the most widely adopted LPWAN standard. Research has focused on Adaptive Data Rate (ADR) algorithms to optimize transmission efficiency and energy use~\cite{lehong2020adr}. Despite these advances, studies consistently highlight LoRaWAN's strict duty cycle and limited payload capacity as major constraints, restricting its use to small telemetry packets rather than large sensor datasets.

In parallel, DTN research has advanced from the original Bundle Protocol (RFC~5050) to the proposed standard BPv7 (RFC~9171)~\cite{burleigh2022bpv7}, which improves efficiency, extensibility, and security. Recent surveys emphasize buffering, custody transfer, and opportunistic forwarding as critical mechanisms for reliability in intermittent, high-latency environments. DTN has now moved into operational deployments, particularly in space communication systems, demonstrating resilience in extreme conditions~\cite{elewaily2024dtn}.

Together, these developments show that current IoT solutions either sacrifice throughput for range (LoRaWAN) or remain confined to specialized contexts (DTN). This gap directly relates to the problem statement: forest environments require both resilience against intermittent connectivity and sufficient bandwidth for richer datasets.

\subsection{Limitations of Existing Solutions}

LoRaWAN, while dominant in LPWAN deployments, suffers from strict duty cycle restrictions and small payload sizes, making it unsuitable for transferring large sensor datasets or multimedia logs. Its reliance on gateway infrastructure also introduces cost and reliability challenges in remote environments.

DTN, though resilient in intermittent and high-latency conditions, faces its own limitations. Despite successful trials, DTN implementations remain tied to specialized hardware or research prototypes, limiting accessibility for commodity deployments~\cite{castillo2024dtn}.

\subsection{Relation to Objectives}

LoRaWAN struggles with payload and duty cycle limits, while DTN faces adoption and complexity challenges. FLP bridges these gaps by combining efficient data transport, delay-tolerant resilience, and adaptive self-healing topologies, tailored for forest connectivity using commodity hardware.

% ===========================================================================
\section{System Design}
\label{sec:design}
% ===========================================================================

\subsection{Architecture Overview}

FLP employs a three-layer communication architecture with adaptive protocol selection, as shown in Fig.~\ref{fig:architecture}:

\begin{enumerate}
  \item \textbf{ESP-NOW (Short-Range Mesh)} --- High-bandwidth (\(\sim\)34\,KB/s per hop), short-range (30--50\,m in forest) inter-node communication for bulk data relay.
  \item \textbf{LoRa (Long-Range Control)} --- Extended range (200--500\,m in forest), lower bandwidth (617\,B/s at SF7) for broadcast discovery, route advertisements, and transfer coordination.
  \item \textbf{WiFi + MQTT (External Gateway)} --- Opportunistic cloud connectivity via MQTT for data exfiltration when a node detects internet access.
\end{enumerate}

\begin{figure}[htbp]
  \centering
  \fbox{\parbox{0.9\linewidth}{\centering\small
    \textbf{FLP Network Topology}\\[4pt]
    ESP32-S3 nodes communicate via ESP-NOW (solid lines) and LoRa (dashed lines).
    Exit nodes with WiFi connectivity forward data to the MQTT broker.
    The cloud-side MQTT Admin reassembles fragments and manages flow control.
  }}
  \caption{FLP multi-protocol mesh architecture. Nodes use ESP-NOW for high-bandwidth data relay and LoRa for broadcast control signaling. Exit nodes bridge to the cloud via MQTT over WiFi.}
  \label{fig:architecture}
\end{figure}

\subsection{Hardware Platform}

All nodes use the LILYGO T3-S3 board, featuring an ESP32-S3 microcontroller with 8\,MB octal PSRAM and an integrated SX1276 LoRa transceiver. Table~\ref{tab:hardware} summarizes the key hardware specifications.

\begin{table}[htbp]
  \centering
  \caption{Hardware Platform Specifications}
  \label{tab:hardware}
  \begin{tabular}{@{}ll@{}}
    \toprule
    \textbf{Component} & \textbf{Specification} \\
    \midrule
    MCU & ESP32-S3, dual-core Xtensa LX7 \\
    PSRAM & 8\,MB octal (fragment reassembly buffer) \\
    LoRa & SX1276, 915\,MHz, SF7--SF10 adaptive \\
    WiFi & 802.11 b/g/n, 2.4\,GHz \\
    ESP-NOW & 250\,B MTU, broadcast + unicast \\
    Display & SSD1306 128$\times$64 OLED \\
    Storage & MicroSD via SPI \\
    \bottomrule
  \end{tabular}
\end{table}

\subsection{Packet Structure}

FLP uses a compact 8-byte packet header (Table~\ref{tab:header}) designed to maximize payload within the 250-byte ESP-NOW MTU and 255-byte LoRa maximum.

\begin{table}[htbp]
  \centering
  \caption{FLP Packet Header (8 bytes)}
  \label{tab:header}
  \begin{tabular}{@{}lll@{}}
    \toprule
    \textbf{Field} & \textbf{Bits} & \textbf{Description} \\
    \midrule
    \texttt{ver\_type} & 8 & Version (2b) + Packet type (6b) \\
    \texttt{src\_addr} & 16 & Source node address \\
    \texttt{dst\_addr} & 16 & Destination address \\
    \texttt{ttl\_hops} & 8 & TTL (4b) + Hop count (4b) \\
    \texttt{seq\_num} & 16 & Sequence number \\
    \bottomrule
  \end{tabular}
\end{table}

Packet types include \texttt{DATA} (0x01), \texttt{ACK} (0x02), \texttt{NACK} (0x03), \texttt{DISCOVERY} (0x10), \texttt{ROUTE\_REQ/REPLY} (0x11/0x12), \texttt{ROUTE\_ERROR} (0x13), \texttt{TRANSFER\_AD} (0x20), \texttt{TRANSFER\_ACK} (0x21), \texttt{PARITY} (0x06), and relay types \texttt{MESH\_PUB/CMD} (0x30/0x31).

\subsection{Routing: Sequence-Numbered Distance Vector}

FLP implements a DSDV-inspired routing protocol with ETX (Expected Transmission Count) weighted metrics. Each discovery broadcast carries:

\begin{itemize}
  \item A monotonic sequence number (\texttt{inet\_seq}) from the originating exit node
  \item The exit node's address (\texttt{inet\_origin})
  \item Hop count to internet
  \item RSSI of the reporting link
  \item WiFi channel and current LoRa spreading factor (packed into a single flags byte)
\end{itemize}

Route updates are accepted only if the sequence number is higher than the current known value (preventing count-to-infinity loops), or equal with a lower composite cost. The composite routing cost is:

\begin{equation}
  \text{cost} = \text{hops\_to\_internet} \times 100 + \text{ETX} \times 100
\end{equation}

where ETX is computed per-neighbor from transmission success rate. This naturally routes around lossy links: a 2-hop path with ETX\,=\,1.0 (cost\,=\,200) beats a 1-hop path with ETX\,=\,5.0 (cost\,=\,600).

\textbf{Route Error Propagation.} When a neighbor becomes unreachable (15\,s without contact), a \texttt{ROUTE\_ERROR} packet is broadcast, enabling upstream nodes to invalidate stale routes within seconds rather than waiting for the 30\,s passive timeout.

\subsection{File Transfer Protocol}

The file transfer process follows a three-phase approach:

\textbf{Phase 1: Exit Node Election.} The source node broadcasts a compact \texttt{TRANSFER\_AD} (9 bytes: session ID, file size, fragment size, CRC-16) via LoRa. Nodes with internet access respond with \texttt{TRANSFER\_ACK} containing their RSSI to gateway and hop count. The source selects up to four exit nodes and assigns fragment ranges round-robin. The election window adapts from 3--8\,s based on the source's hop distance to the nearest known exit.

\textbf{Phase 2: Pipelined Data Transfer.} Data is fragmented into 240-byte chunks and transmitted using Selective Repeat ARQ with a sliding window of 32 fragments per exit node. Fragment~0 carries a filename prefix (\texttt{[len:1][filename:N][data...]}) so the exit node can publish transfer metadata to MQTT after receiving it. Every 7th fragment slot carries a FEC parity fragment (XOR of the preceding group), enabling recovery from single fragment losses without retransmission.

\textbf{Phase 3: Cloud Reassembly.} Exit nodes forward fragments to the MQTT broker as \texttt{[2-byte seq][data]} payloads on the topic \texttt{flp/<exit\_addr>/file/data}. The cloud-side MQTT Admin reconstructs the file from fragments, using session ID correlation and CRC verification.

\subsection{Adaptive Protocol Selection}

A \texttt{ProtocolSelector} module continuously evaluates transport suitability based on:

\begin{itemize}
  \item Packet size (LoRa favored for small control packets)
  \item ESP-NOW neighbor reachability
  \item Historical link success rate per transport
  \item Hop count (LoRa for broadcast discovery, ESP-NOW for unicast data)
\end{itemize}

\subsection{Adaptive LoRa Spreading Factor}

Inspired by LoRaWAN's ADR mechanism, FLP adjusts the LoRa spreading factor (SF7--SF10) every 10\,s based on link quality and neighbor count. Poor LoRa success rates with few neighbors trigger an increase to SF10 for extended range, while good conditions with sufficient neighbors revert to SF7 for maximum throughput.

\subsection{Security}

FLP employs hardware-accelerated AES-128-GCM encryption on the ESP32-S3, providing both confidentiality and integrity for all mesh traffic without significant throughput reduction.

\subsection{Data Storage and Processing}

Collected data is stored in the ESP32-S3's 8\,MB PSRAM until internet availability is detected. Fragment reassembly uses $O(1)$ memory-mapped indexing: each fragment's offset is computed directly from its sequence number, enabling reassembly without sorting or searching. Once published to the cloud, the MQTT Admin reconstructs fragmented packets, discards duplicates, and rebuilds the original file.

\subsection{Cloud Interaction}

Users interact with forest nodes through the cloud by publishing MQTT packets to the MQTT Admin, which are relayed back down through the mesh to specific nodes via the \texttt{MESH\_CMD} packet type. Each node maintains a synchronized topic lookup table for efficient relay-topic mapping.

% ===========================================================================
\section{Implementation}
\label{sec:implementation}
% ===========================================================================

\subsection{System Prototyping and Development Tools}

The FLP system is prototyped using ESP32-S3 microcontrollers equipped with 8\,MB PSRAM, programmed in C/C++ using ESP-IDF v5.5 (Espressif IoT Development Framework). ESP-IDF provides the build system, component management, and runtime configuration through Kconfig. The dual-protocol mesh relies on ESP-NOW for short-range, high-bandwidth communication and an SX1276 LoRa transceiver on the LILYGO T3-S3 board for long-range links.

On the cloud side, a Mosquitto MQTT broker serves as the external gateway, paired with an MQTT Admin service responsible for fragment reassembly, sliding window flow control, and topic management. Each node is fitted with an SSD1306 OLED display for real-time status output during development and demonstration.

Key software tools include:
\begin{itemize}
  \item ESP-IDF v5.5 with FreeRTOS for multi-threaded task management
  \item VS Code with clangd for Barr-C coding standards enforcement
  \item Python 3.8+ for the ESP-IDF toolchain
\end{itemize}

The firmware is organized into independent ESP-IDF components: \texttt{espnow/} (ESP-NOW transport), \texttt{lora/} (LoRa transport), \texttt{mesh/} (MeshManager, TransferEngine, route table), \texttt{protocol/} (ARQ, FEC, packet definitions), \texttt{mqtt\_client/} (MQTT-SN client), \texttt{display/} (OLED driver), and \texttt{sdcard/} (SD card reader). Each component is developed and validated independently before integration.

\subsection{Deployment Strategy}

The FLP system is deployed in a multi-layered mesh configuration consisting of Sensor Nodes (origin nodes), Relay Nodes (extenders), and Cloud Gateways (data exit points). Deployment in a dense forest presents the ``Green Wall'' effect where moisture-heavy foliage attenuates 2.4\,GHz signals. FLP mitigates this through:

\begin{enumerate}
  \item \textbf{Adaptive Protocol Selection} --- FLP uses LoRa as a control plane to discover other nodes. Once a neighbor is identified, FLP escalates to ESP-NOW for high-speed file transfers. If the ESP-NOW RSSI drops below a threshold, FLP falls back to LoRa.
  \item \textbf{Network Resilience} --- To prevent management collapse observed in TCP/IP meshes, FLP uses Opportunistic Exit Node Election, where nodes with internet access advertise their state and sensor nodes broadcast an ``intent'' packet via LoRa to locate these exit nodes dynamically.
  \item \textbf{Priority Scheduling} --- A dual-priority FreeRTOS queue ensures that control packets (ACK, NACK, route updates) are processed before bulk data fragments, preventing ARQ timeouts under heavy transfer load.
\end{enumerate}

\subsection{Testing Process}

For this project, the system is deployed and tested in a lab environment. Nodes are placed at intervals across the lab space, with at least one node within WiFi range of an internet-connected access point to serve as an exit node. The remaining nodes communicate exclusively over ESP-NOW and LoRa, relaying data hop-by-hop toward the nearest exit node. Physical obstructions such as walls, furniture, and partitions within the lab are used to simulate signal attenuation and approximate the NLOS conditions of a real forest deployment.

\subsection{Performance Benchmarks}

Success is evaluated against the following metrics:
\begin{itemize}
  \item \textbf{NFR-MESH1}: Deliver a 1\,MB file over 3+ hops within 20 minutes
  \item \textbf{End-to-end throughput}: Measured in KB/s across varying hop counts
  \item \textbf{Packet delivery ratio}: Percentage of fragments successfully delivered
  \item \textbf{Route convergence time}: Time for the network to stabilize after topology changes
  \item \textbf{ARQ retransmission rate}: Indicator of link quality and FEC effectiveness
\end{itemize}

% ===========================================================================
\section{Results and Analysis}
\label{sec:results}
% ===========================================================================

\subsection{Throughput Analysis}

Table~\ref{tab:throughput} presents the per-component throughput analysis before and after optimization of the FLP stack.

\begin{table}[htbp]
  \centering
  \caption{Per-Component Throughput: Before vs After Optimization}
  \label{tab:throughput}
  \begin{tabular}{@{}lrrl@{}}
    \toprule
    \textbf{Component} & \textbf{Before} & \textbf{After} & \textbf{Gain} \\
    \midrule
    ESP-NOW effective throughput & 8\,KB/s & 34\,KB/s & 4.3$\times$ \\
    ARQ pipeline capacity & 6.6\,KB/s & 79\,KB/s & 12$\times$ \\
    MQTT drain rate & 2\,frag/s & 200+\,frag/s & 100$\times$ \\
    Fragment queue buffer & 0.1\,s & 1.5\,s & 15$\times$ \\
    \bottomrule
  \end{tabular}
\end{table}

The primary bottleneck shifted from software limitations (MQTT drain delay, small ARQ window) to physics (radio air-time sharing at relay nodes), indicating the software stack is no longer the limiting factor.

Key optimizations included: enlarging the ARQ window from 8 to 32 fragments, reducing the MQTT processing interval from 500\,ms to 10\,ms, expanding the fragment queue from 8 to 64 entries, and enabling BLE Data Length Extension for larger link-layer PDUs.

\subsection{Transfer Time Performance}

Table~\ref{tab:transfer_time} shows the projected 3\,MB file transfer times across different environmental conditions with four exit nodes.

\begin{table}[htbp]
  \centering
  \caption{3\,MB Transfer Time Over 3 Hops (4 Exit Nodes)}
  \label{tab:transfer_time}
  \begin{tabular}{@{}lr@{}}
    \toprule
    \textbf{Condition} & \textbf{Time} \\
    \midrule
    Lab (clean RF, short range) & $\sim$50\,s \\
    Moderate forest (partial canopy) & $\sim$70\,s \\
    Dense forest (full canopy, dry) & $\sim$85\,s \\
    Dense forest + rain & $\sim$95\,s \\
    Worst case (heavy rain, high humidity) & $\sim$110\,s \\
    \bottomrule
  \end{tabular}
\end{table}

All scenarios meet the NFR-MESH1 target of delivering data within 20 minutes, with dense forest conditions achieving transfer times of approximately 1.5--2 minutes---well within the requirement.

\subsection{Scalability}

Table~\ref{tab:scalability} presents the relationship between network size and transfer performance in a Singapore tropical forest environment.

\begin{table}[htbp]
  \centering
  \caption{Performance Scaling with Network Size (Singapore Forest)}
  \label{tab:scalability}
  \begin{tabular}{@{}ccrrc@{}}
    \toprule
    \textbf{Nodes} & \textbf{Hops} & \textbf{LoRa Range} & \textbf{3\,MB Time} & \textbf{Reliability} \\
    \midrule
    2 & 1 & 200\,m & 2.0\,min & $>$98\% \\
    4 & 3 & 600\,m & 2.6\,min & $>$95\% \\
    6 & 5 & 1.0\,km & 2.8\,min & $\sim$91\% \\
    9 & 8 & 1.6\,km & 3.0\,min & $\sim$82\% \\
    \bottomrule
  \end{tabular}
\end{table}

Transfer time scales slowly with hop count (2.0 to 3.0 minutes across 1--8 hops) because the pipelined ARQ keeps all relay links busy in parallel. The relay node's per-hop throughput ($\sim$25\,KB/s), not the number of hops, is the dominant bottleneck.

\subsection{Routing Performance}

The DSDV sequence-numbered routing with ETX weighting provides several measurable improvements over the baseline hop-count-only metric:

\begin{itemize}
  \item \textbf{Route convergence}: Stale routes are invalidated within seconds via \texttt{ROUTE\_ERROR} propagation, compared to the 30\,s passive timeout in the baseline.
  \item \textbf{Lossy link avoidance}: ETX-weighted routing naturally avoids high-loss links, reducing the overall ARQ retransmission rate.
  \item \textbf{Congestion-aware forwarding}: When the packet queue exceeds 75\% capacity, low-priority forwarded packets are dropped, protecting control traffic from being starved during heavy transfers.
\end{itemize}

\subsection{Comparison with Existing Solutions}

Table~\ref{tab:comparison} compares FLP against LoRaWAN and DTN for the target use case of multi-megabyte forest data transfer.

\begin{table}[htbp]
  \centering
  \caption{Comparison with Existing Protocols}
  \label{tab:comparison}
  \begin{tabular}{@{}lccc@{}}
    \toprule
    \textbf{Metric} & \textbf{FLP} & \textbf{LoRaWAN} & \textbf{DTN} \\
    \midrule
    3\,MB transfer & $\sim$2\,min & $>$85\,min & Variable \\
    Max payload/pkt & 242\,B & 51\,B & Unlimited \\
    Self-healing & Yes & No (star) & Partial \\
    Hardware cost & \$15/node & \$30+ GW & Specialized \\
    Duty cycle limit & None & 1\% & None \\
    \bottomrule
  \end{tabular}
\end{table}

\subsection{Bottlenecks and Limitations}

The primary remaining bottleneck is ESP-NOW radio air-time sharing at relay nodes. Since each ESP32-S3 has a single 2.4\,GHz radio that must be time-shared between receiving from upstream and transmitting to downstream, the effective per-hop throughput is halved at relay nodes ($\sim$34\,KB/s effective vs $\sim$81\,KB/s raw). This is a fundamental hardware constraint rather than a protocol limitation.

The LoRa control plane uses unslotted ALOHA, limiting the maximum number of nodes in a single hearing area to approximately 14 before contention degrades performance. Larger deployments require network segmentation into clusters.

% ===========================================================================
\section{Conclusion}
\label{sec:conclusion}
% ===========================================================================

This paper presented Forest Link Protocol, a multi-protocol adaptive mesh network designed for reliable multi-megabyte data transfer in dense forest environments. FLP addresses the Bandwidth-Reach Paradox by combining ESP-NOW for high-throughput short-range relay with LoRa for long-range control signaling, bridging the gap between LoRaWAN's payload constraints and traditional WiFi mesh fragility.

Our implementation on commodity ESP32-S3 hardware demonstrates that a 3\,MB file can be transferred over three hops in approximately 1.5--2 minutes under dense forest conditions, exceeding the 20-minute target requirement by an order of magnitude. The DSDV-based routing with ETX weighting provides rapid route convergence and automatic avoidance of lossy links, while the dual-priority packet scheduling ensures control traffic is not starved during heavy data transfers.

FLP contributes to the broader field of IoT by demonstrating that commodity hardware, optimized at the software layer, can perform comparably to specialized LPWAN infrastructure at a fraction of the cost. The protocol's self-healing topology and delay-tolerant design make it suitable for real-world deployments in environmental monitoring, wildfire detection, and biodiversity data collection.

Future work includes: (1) implementing LoRa TDMA scheduling to reduce control-plane collisions beyond 14 nodes, (2) field trials in Singapore's Central Catchment Nature Reserve with 15--25 nodes, (3) integrating BLE 5.0 Coded PHY for extended control signaling range, and (4) evaluating energy harvesting integration for truly autonomous long-term deployments.

% ===========================================================================
% References
% ===========================================================================
\balance
\begin{thebibliography}{5}

\bibitem{maurya2025lorawan}
P.~Maurya, A.~Hazra, P.~Kumari, T.~B.~S{\o}rensen, and S.~K.~Das,
``A comprehensive survey of data-driven solutions for LoRaWAN: Challenges and future directions,''
\textit{ACM Trans.\ Internet Things}, vol.~6, no.~1, pp.~1--36, Feb.\ 2025,
doi: 10.1145/3711953.

\bibitem{elewaily2024dtn}
D.~I.~Elewaily, H.~A.~Ali, A.~I.~Saleh, and M.~M.~Abdelsalam,
``Delay/disruption-tolerant networking-based the integrated deep-space relay network: State-of-the-art,''
\textit{Ad Hoc Networks}, vol.~152, pp.~103307, Jan.\ 2024,
doi: 10.1016/j.adhoc.2023.103307.

\bibitem{lehong2020adr}
C.~Lehong, B.~Isong, F.~Lugayizi, and A.~M.~Abu-Mahfouz,
``A survey of LoRaWAN adaptive data rate algorithms for possible optimization,''
in \textit{Proc.\ 2nd Int.\ Multidisciplinary Inf.\ Technol.\ Eng.\ Conf.\ (IMITEC)},
pp.~1--9, Nov.\ 2020,
doi: 10.1109/imitec50163.2020.9334144.

\bibitem{burleigh2022bpv7}
S.~Burleigh, K.~Fall, and E.~Birrane,
``Bundle Protocol Version 7,''
RFC~9171, Feb.\ 2022,
doi: 10.17487/rfc9171.

\bibitem{castillo2024dtn}
A.~Castillo, C.~Juiz, and B.~Bermejo,
``Delay and disruption tolerant networking for terrestrial and TCP/IP applications: A systematic literature review,''
\textit{Network}, vol.~4, no.~3, pp.~237--259, Jul.\ 2024,
doi: 10.3390/network4030012.

\end{thebibliography}

\end{document}
